\begin{theorem}[Failure of global weak uniform continuity]\label{uniform-failure}
On the balanced unit-weight line with actual reference $\gamma(x)=|x|$,
the finite raw operator is not globally uniformly continuous for the weak
uniformity generated by $C_b(\R)$, even on literal relative PA PC functions.
On the real relative PA domain it is not globally uniformly continuous
for the bounded-Lipschitz metric either.
\end{theorem}
\begin{proof}
For integers $n\geq1$, put $A_n=|x-2n|$ and
$B_n=|x-2n-1/(8n)|$. A carrying structure has the rational vertices
$0,2n,2n+1/(8n)$ and their complementary edges. The affine formulas
have slopes $\pm1$ and rational constants. Both functions are ambient
convex, so satisfy every printed rational-endpoint PCC. Their differences
from the actual $|x|$ are bounded, with norms $2n$ and $2n+1/(8n)$;
their two recession slopes are one. They are actual literal PA PC
endpoints with constant decreasing regularizations. Primitive outgoing
normals on the line are $\pm1$. Away from each kink the opposite slopes
cancel, while the two outgoing slopes at the kink are both one. Thus
\[
 \MA(A_n)=2\delta_{2n},\quad \MA(B_n)=2\delta_{2n+1/(8n)},\quad
 \|A_n-B_n\|_\infty=1/(8n).
\]
For the single fixed test $\psi(x)=\sin(\pi x^2)$, which is bounded
and continuous, the absolute evaluation difference is
$2\cos(\pi/(64n^2))\to2$. This violates one fixed defining entourage
of the weak $C_b$ uniformity.

For even $n$ let $a_j=4j$ for even $j$, $a_j=4j+1$ for odd $j$,
and let $b_j$ make the opposite choices, $0\leq j<n$. Set
$A_n=n^{-1}\sum_j|x-a_j|$, $B_n=n^{-1}\sum_j|x-b_j|$.
Each sum is a convex function on the rational line subdivision by these
finitely many integer points, with rational slopes and constants. Its
outer slopes are $-1,1$; hence it has bounded difference from $|x|$
and the prescribed positive recession growth. It belongs to real PA on
finite rational carrying geometry and is already regularizable there.
This need not be literal PA. Raw linearity in the actual one-cycle gives
\[
 \MA(A_n)=\frac2n\sum_j\delta_{a_j},\qquad
 \MA(B_n)=\frac2n\sum_j\delta_{b_j}.
\]
For a pair $j=2k,2k+1$, its unnormalized difference is
\[
 |x-8k|-|x-8k-1|+|x-8k-5|-|x-8k-4|.
\]
It vanishes outside $[8k,8k+5]$, rises linearly from zero to two
on $[8k,8k+1]$, equals two on $[8k+1,8k+4]$, and falls to zero
on $[8k+4,8k+5]$. These pair intervals are disjoint. Therefore
$0\leq A_n-B_n\leq2/n$ and the uniform distance is exactly $2/n$.
Let $\eta$ be the continuous eight-periodic linear interpolation through
$(0,1/4),(1,-1/4),(4,-1/4),(5,1/4),(8,1/4)$.
It has sup norm $1/4$ and Lipschitz constant $1/2$; it belongs even
to the convention $\|\eta\|_\infty+\operatorname{Lip}(\eta)\leq1$.
Its values at every $a_j,b_j$ are $1/4,-1/4$, respectively, so the
raw measure evaluation difference is exactly one for every even $n$.

The relative norm of a sum of these shifted absolute values against
$|x|$ is exactly its average shift: the triangle inequality gives the
upper bound, attained on $x\leq0$. Both averages equal
$2(n-1)+1/2=2n-3/2$. They grow without a common fixed bound. These
families are not uniformly convergent to one relative function. The
failure concerns global uniform continuity, not the convergence result
proved next. It is a material restriction of the old guarantee.
\end{proof}

We specify all output notions. The vague topology and uniformity are
generated by finite lists of evaluations $\nu\mapsto\int g\,d\nu$,
$g\in C_c(X)$. The weak $C_b$ topology and uniformity use finite lists
with $g\in C_b(X)$. Fix a Euclidean metric on $X$ and define
\[
 d_{\rm BL}(\mu,\nu)=\sup_{\|g\|_\infty\leq1,\ \operatorname{Lip}(g)\leq1}
              \left|\int g\,d(\mu-\nu)\right|.
\]
All measures below have the same finite raw mass $M$. The domain metric
on tuples is $d_\infty(\mathbf f,\mathbf g)=\sum_i\|f_i-g_i\|_\infty$.
For each category $*$, let $\mathcal D_*(\gamma)$ be the uniform closure
of the actual relative $\PAPC_*$ functions in continuous functions of
bounded difference from the same actual reference. No other continuous
PC function is added by definition.

\begin{lemma}[Positive-functional representation]\label{radon-lemma}
On the closed finite polyhedral metric space $X$, every positive linear
functional $L$ on $C_c(X)$ satisfying $|L(g)|\leq M\|g\|_\infty$ is
integration against a unique finite positive Radon measure of mass at
most $M$.
\end{lemma}
\begin{proof}
For open $O$, define
$\alpha(O)=\sup\{L(h):h\in C_c(X),0\leq h\leq1,
\operatorname{supp}h\subset O\}$.
It is at most $M$. If $O\subset\bigcup_iO_i$, the compact support of
any such $h$ has a finite subcover. A continuous partition on a neighborhood
of that support, with compact supports subordinate to the selected opens,
exists by taking finitely many small relatively compact balls, continuous
distance cutoffs, and dividing their sum where it is positive. It splits
$h$ into nonnegative functions $h_i\leq1$ supported in $O_i$.
Hence $\alpha(O)\leq\sum_i\alpha(O_i)$. For disjoint opens finite
sums of admissible tests are again at most one; taking suprema gives
$\alpha(\bigcup_iO_i)\geq\sum_i\alpha(O_i)$, first for finite sums
and then for all finite partial sums.

Set $\mu^*(A)=\inf_{O\supset A}\alpha(O)$ over open supersets.
Choosing open supersets within arbitrary summable errors proves countable
subadditivity, so this is an outer measure. If two sets have positive
distance, their disjoint distance neighborhoods, intersected with any
open superset of their union, give
$\mu^*(A\cup B)=\mu^*(A)+\mu^*(B)$.
Here is the Borel measurability step in detail. For a closed $E$ and
arbitrary $A$, put $A_j=A\cap\{\operatorname{dist}(x,E)\geq2^{-j}\}$.
Separated-set additivity gives
$\mu^*(A)\geq\mu^*(A\cap E)+\mu^*(A_j)$.
The shells of $A\setminus E$ between successive positive dyadic distances
have finite total outer mass: any finite collection of even-index shells
is pairwise separated and has summed outer mass at most $M$, and the
same is true for odd-index shells. Consequently the sums of all shell
masses converge and their tails tend to zero. Countable subadditivity
then gives
$\mu^*(A\setminus E)\leq\lim_j\mu^*(A_j)$.
This is the Caratheodory inequality for $E$; its reverse is subadditivity.
Thus all closed, hence all Borel, sets are measurable. The restriction
$\mu$ is a countably additive finite measure, since the outer-measure
identity for measurable disjoint sets and continuity of finite partial
sums give countable additivity.

For open $O$ one has $\mu(O)=\alpha(O)$ directly by monotonicity in
open supersets. If $h$ is an admissible test and $K=\operatorname{supp}h$,
then $L(h)\leq\alpha(W)$ for every open $W\supset K$, whence
$L(h)\leq\mu(K)$. This proves inner compact regularity on opens;
outer regularity is the definition of $\mu^*$. For a Borel set, choose
an open superset of arbitrarily small excess measure and apply inner
regularity to its complement in a sufficiently large compact set.
Since $\alpha(X)$ is approximated by compactly supported tests, such
compact sets exhaust its mass. This proves inner regularity on all
Borel sets as well.

For $g\in C_c(X)$ nonnegative, let $R=\|g\|_\infty$ and $\delta=R/n$.
Tests $h_k$ supported in $\{g>k\delta\}$ satisfy
$\delta\sum_{k=1}^n h_k\leq g$; maximizing each of their functional
values gives the lower bound
$L(g)\geq\delta\sum_{k=1}^n\mu\{g>k\delta\}$.
For the upper bound choose a compact cutoff equal to one on
$\operatorname{supp}g$, and for each $k\geq1$ a test equal to one
on the compact set $\{g\geq k\delta\}$ and supported in
$\{g>(k-1)\delta\}$. Their sum times $\delta$, together with
$\delta$ times the first cutoff, majorizes $g$.
Thus
$L(g)\leq\delta M+\delta\sum_{k=1}^n\mu\{g>(k-1)\delta\}$.
The two sums differ by at most $\delta M$ and also bracket the integral
of $g$ by elementary pointwise step bounds. Letting $n\to\infty$
proves $L(g)=\int g\,d\mu$. Signed functions follow by their positive
and negative parts. Conversely, compact cutoffs and inner regularity
recover open-set masses from these integrals, proving uniqueness.
\end{proof}

\begin{theorem}[Raw completion with the exact continuity scopes]\label{raw-extension}
Fix finite pure positively balanced embedded rational $X$ and the actual
rational-affine PA PC positive-growth reference $\gamma$ of this section.
For $*=\Z,\R$, the finite raw mixed operator extends uniquely to
$\mathcal D_*(\gamma)^d$. Its outputs are positive finite Radon measures
of raw mass $M=\deg(\gamma)$; the extension is symmetric, local, and
agrees with the finite raw recursion. Arbitrary uniformly convergent PA
approximants give the same limit against every $C_b$ test.

The extension is continuous from the uniform domain topology into the
weak $C_b$ topology (for sequences and nets). It is globally uniformly
continuous into the vague $C_c$-test uniformity. For \emph{each fixed}
$0\leq C<\infty$, on
\[
 \mathcal F_{*,C}=\{(f_1,\ldots,f_d)\in\mathcal D_*(\gamma)^d:
                   \max_i\|f_i-\gamma\|_\infty\leq C\},
\]
it is uniformly continuous into the weak $C_b$ uniformity and in
$d_{\rm BL}$. These last two assertions have no entire-domain
uniform-continuity conclusion. If $M=0$, all outputs are zero and
all the asserted continuity properties are immediate.
\end{theorem}
\begin{proof}
Assume $M>0$. For a compact ambient rational PA test $\psi$, the preceding
DC lemma gives $\theta_1-\theta_0=\varepsilon\psi$ with actual relative
real PA PC functions and fixed $\varepsilon>0$. For two relative finite
PA PC tuples, telescope one factor at a time. The other factors form
positive balanced one-cycles $Z_i$, so bounded edge integration by parts
and raw positivity give
\begin{equation}\label{test-estimate}
 \left|\int\psi\bigl(\MA(\mathbf a)-\MA(\mathbf b)\bigr)\right|
 \leq \frac{2M}{\varepsilon}\sum_i\|a_i-b_i\|_\infty.
\end{equation}
Indeed $D_\psi Z_i=(D_{\theta_1}Z_i-D_{\theta_0}Z_i)/\varepsilon$
has total variation at most $2M/\varepsilon$. This constant depends on
the test and actual reference, \emph{not} on relative sizes of $\mathbf a$
or $\mathbf b$. Auxiliary real functions are legitimate even when the
approximating tuples are literal, because their finite real raw recursion
has already been established; no new literal endpoint is inferred.

For a uniform-Cauchy tuple of approximants, the integrals in
\eqref{test-estimate} are Cauchy. Uniform density of compact rational PA
tests in $C_c(X)$, with supports controlled as proved above, and the
common mass $M$ show that every $C_c$ integral is Cauchy: replacing a test
within $\eta$ costs at most $2M\eta$ in a difference. The limits define
a positive linear functional bounded by $M\|g\|_\infty$.
Lemma~\ref{radon-lemma} gives its unique positive Radon measure $\mu$.

We prove mass and tightness without changing the actual reference.
Choose a fixed rational $c\geq0$ so $H=\gamma+c\geq0$ on $X$.
This is possible by the coercive lower bound; $H$ is a proper height.
The raw reference measure remains $\MA(\gamma)$ by constant invariance.
Put $J=\int H\MA(\gamma)<\infty$, and for integers $j\geq1$ set
\[
 q_j=\max(\gamma+c-j,0)-c=\max(\gamma-j,-c).
\]
These are relative rational-affine PA PC functions on finite rational
structures; $\gamma-q_j=\min(H,j)$.
For a finite tuple with $\max_i\|a_i-\gamma\|_\infty\leq D$,
apply the previously proved relative IBP identity with $q=q_j$.
Split $a_1-q_j=(a_1-\gamma)+(\gamma-q_j)$.
The first two terms have absolute values at most $DM$ each, by positive
mixed measures of mass $M$. The remaining term replaces $a_1$ by
$\gamma$. Repeat through all $d$ factors, obtaining
\begin{equation}\label{moment}
 \int\min(H,j)\MA(\mathbf a)\leq2dDM+J=:K_D.
\end{equation}
Monotone convergence over $j$ gives the same bound for $\int H$.
Consequently $\MA(\mathbf a)\{H>T\}\leq K_D/T$.

A particular uniformly convergent family of approximants has one finite
$D$. Choose continuous compact cutoffs $\chi_T$ equal to one on
$\{H\leq T\}$, with $0\leq\chi_T\leq1$.
Vague convergence gives $\int\chi_T\,d\mu\geq M-K_D/T$ and the
representation gives $\mu(X)\leq M$. Hence $\mu(X)=M$.
Nonnegative compact tests majorized by $H$ pass to the limit with the
same bound $K_D$. Their suprema, or an increasing exhaustion of $H$ by
compact cutoffs, give $\int H\,d\mu\leq K_D$.
Truncating a fixed $g\in C_b(X)$ by $\chi_T$ costs at most
$\|g\|_\infty K_D/T$ for either the approximant or the limit.
Vague convergence of the truncated test, then $T\to\infty$, proves
convergence against \emph{each} $C_b$ test for this uniformly convergent
family. This reasoning has not bounded every relative tuple by one $D$.

Comparing any two approximating tuples by \eqref{test-estimate} proves
sequence independence. Passing that estimate along approximants for
any two completed tuples gives \emph{the same} compact-test estimate
on the completion. If $\mathbf f\in\mathcal F_{*,C}$, choose approximants
with uniform error tending to zero, so their relative bounds are
$C+o(1)$. Passing each nonnegative compact height test in
\eqref{moment} and then taking its supremum gives the exact bound
\begin{equation}\label{completed-moment}
 \int H\,d\MA(\mathbf f)\leq K_C=2dCM+J,
 \qquad \MA(\mathbf f)\{H>T\}\leq K_C/T
       \quad(\mathbf f\in\mathcal F_{*,C}).
\end{equation}
Thus the common constant applies to \emph{all} completed tuples in
this fixed-$C$ family, not just a selected sequence. Using approximants
within one would also give the valid weaker constant $K_{C+1}$.

For a fixed $g\in C_c(X)$ approximate it uniformly within $\eta$ by
one compact rational PA test $\psi$. For \emph{all} completed tuples,
\[
 |\langle g,\MA(\mathbf f)-\MA(\mathbf h)\rangle|
 \leq2M\eta+L_\psi d_\infty(\mathbf f,\mathbf h),\qquad
 L_\psi=2M/\varepsilon_\psi.
\]
First choose $\eta$, then the input distance. Finite lists use the
minimum of finitely many input thresholds. This proves global uniform
continuity for the \emph{vague} uniformity.

Now fix $C$ once and for all and $g\in C_b(X)$. With the preceding
cutoff and a compact rational PA approximation to $\chi_Tg$ within
$\eta$, for every $\mathbf f,\mathbf h\in\mathcal F_{*,C}$ we have
\begin{equation}\label{weak-fixed-C}
 |\langle g,\MA(\mathbf f)-\MA(\mathbf h)\rangle|
 \leq2\|g\|_\infty K_C/T+2M\eta+
                          L_\psi d_\infty(\mathbf f,\mathbf h).
\end{equation}
Choose $T,\eta$, then distance. Each is independent of the two tuples
in this fixed family. The same finite-list argument proves uniform
continuity for the weak $C_b$ uniformity on $\mathcal F_{*,C}$.
At one completed tuple $\mathbf f$, its input ball of radius one is
contained in $\mathcal F_{*,C_f}$ where
$C_f=1+\max_i\|f_i-\gamma\|_\infty$.
The last assertion proves ordinary continuity at $\mathbf f$ into
weak $C_b$, including nets. A uniformly convergent family has a common
finite $C$; this is the exact scope of the all-$C_b$ convergence claim.

For the bounded-Lipschitz metric, use the same fixed $C$ and cutoff
$\chi_T$, and let $K=\operatorname{supp}\chi_T$.
The restrictions to $K$ of tests $\|g\|_\infty\leq1$,
$\operatorname{Lip}(g)\leq1$ have a finite uniform $\eta$-net.
To prove this, cover $K$ by finitely many balls of radius $r$ and
quantize the values at their centers with mesh $s$. Two tests in the
same occupied quantization class differ on $K$ by at most $2r+2s$.
Choose one actual global test from each nonempty class; there are finitely
many classes, and choosing $r,s$ small gives the asserted net.
Let its representatives be $g_1,\ldots,g_l$ and approximate each
$\chi_Tg_k$ by a compact rational PA $\psi_k$ within $\eta$.
For any test $g$, choosing the matching representative bounds the
evaluation difference for all tuple pairs in $\mathcal F_{*,C}$ by
\[
 2K_C/T+4M\eta+
       \bigl(\max_k L_{\psi_k}\bigr)d_\infty(\mathbf f,\mathbf h).
\]
All constants are finite and tuple-independent on the \emph{fixed-$C$}
family. Taking the supremum over $g$ proves $d_{\rm BL}$ uniform
continuity there. The counterexamples in Theorem~\ref{uniform-failure}
prevent suppressing that quantifier.

Symmetry follows from finite raw symmetry. For locality, let two limiting
tuples agree in an open $O$, and take a compact rational PA test supported
inside $O$. In the telescoping proof $D_\psi Z_i$ is supported in that
compact subset of $O$, since the test is identically zero on its open
complement. Its variation bound therefore multiplies only the argument
difference on that support. Along approximants to the two tuples this
local difference tends to zero. Density of compact tests with supports
inside $O$ proves equality of output restrictions there. Controlled
support density follows by choosing the bump radii smaller than the
positive distance of the compact test support to $X\setminus O$.
Finite constant sequences show agreement with the original raw operator;
positivity has been preserved at every construction step.

Finally choose approximants $a_n$ with $\|a_n-f\|_\infty\leq2^{-n}$
and set $b_n=a_n+4\cdot2^{-n}$. Then
$b_n-b_{n+1}\geq2\cdot2^{-n}-(2^{-n}+2^{-n-1})=2^{-n-1}>0$
everywhere, and $b_n\to f$ uniformly. These rational additions preserve
both coefficient categories and relative growth. Thus uniform completion
is exactly the declared decreasing regularizable domain. Uniqueness
follows from density of finite approximants and the Hausdorff vague
measure topology. No operator is assigned merely from PC continuity.
\end{proof}

\begin{corollary}[Stability of domain-checked contact integrals]\label{contact-stability}
For the actual reference of Theorem~\ref{raw-extension}, let $u_n,u\in C_b(X)$,
$u_n\to u$ uniformly, and assume the actual envelopes $f_n,f$ belong to
$\mathcal D_*(\gamma)$. Then their exact contact integrals converge.
One may fix any finite $C\geq\sup_n\|u_n\|_\infty,\|u\|_\infty$;
all these envelopes lie in $\mathcal F_{*,C}$ and their measures converge
in the weak $C_b$ topology. This is a convergence theorem with these
domain and common-bound antecedents, not an all-obstacle regularity theorem.
\end{corollary}
\begin{proof}
R230 applies on the actual finite positive balanced space with continuous
PC reference and bounded continuous obstacles. It gives continuous
unrestricted envelopes, the sandwich $\gamma-\|u_n\|_\infty\leq f_n
\leq\gamma+\|u_n\|_\infty$, and $\|f_n-f\|_\infty\leq\|u_n-u\|_\infty$.
Their bounded continuous gaps $f_n-\gamma-u_n$ converge uniformly.
The fixed-$C$ weak convergence theorem and the common mass $M$ then
prove convergence of the gap integrals. Domain membership was assumed,
not deduced from this argument.
\end{proof}
